We evaluated our three dropout checkpoints on the validation pool (CPU, dropout off) and averaged probabilities with each other and with the supplied no-dropout members (Table~\ref{tab:comp}). One caveat up front: our dropout members all use seed 42 and the same data order, so they differ only through the dropout rate, whereas the supplied bank differs by seed.
\begin{itemize}
\item Dropout and ensembling are partial substitutes. The single-model gain of dropout .05 is 0.0036; paired with one supplied member it keeps 0.0011 of that (1.2564 against 1.2576 for two no-dropout seeds), and at $K=5$ swapping one bank member for the dropout .05 model gives 1.2540 against 1.2543. Most of what dropout buys, averaging already buys.
\item Members must be good to help. Adding the .15 and .25 models to the .05 model gives 1.2573 and then 1.2574 (no change for the third), while three no-dropout seeds reach 1.2557: a weaker member dilutes the average, and same-seed members add little diversity. Filling the ensemble with bank members after the three dropout models (1.2557 at $K=5$) is still 0.0014 behind the all-bank ensemble.
\item Rule: build ensembles from the best independently trained members (different seeds, and light dropout if it helps the single model); do not trade seed diversity for regularization diversity.
\end{itemize}
